% Editable conceptual figure for the combined physical/latent diffusion paper.
% Render with: python3 make_conceptual_figure.py
% All graphs and images below are schematic; no checkpoint or data is read.
% Six visible nodes and two Ising nodes represent subsets of the full network.
% Additional per-pixel Ising units, present in the selected Potts1536 run, are omitted.
% The Potts box is ONE unit with mutually exclusive states, not M independent units.
% Local learning denotes the score-matching gradient, including current weights
% and node-level aggregate signals. It is not a weight-independent Hebbian rule.
% The sampling equation describes the effective, fast-hidden-equilibration limit.
% No simulated finite-speed hidden dynamics or fabricated hardware is implied.
\documentclass{article}
\usepackage[paperwidth=19cm,paperheight=13.2cm,margin=0cm]{geometry}
\usepackage[T1]{fontenc}
\usepackage{lmodern}
\usepackage{amsmath,amssymb}
\usepackage{tikz}
\usetikzlibrary{arrows.meta,calc,decorations.pathmorphing}
\renewcommand{\familydefault}{\sfdefault}
\pagestyle{empty}
\definecolor{ink}{HTML}{263442}
\definecolor{muted}{HTML}{596977}
\definecolor{edgeblue}{HTML}{0072B2}
\definecolor{hiddenfill}{HTML}{E5EFF6}
\definecolor{visiblefill}{HTML}{F6F8FA}
\definecolor{springgray}{HTML}{8A959E}
\definecolor{rulegray}{HTML}{D9E0E5}
\definecolor{controlfill}{HTML}{EEF3F7}
\definecolor{accent}{HTML}{D55E00}
\tikzset{
  every node/.style={text=ink,font=\fontsize{8.5}{10.3}\selectfont},
  paneltitle/.style={anchor=west,font=\fontsize{10}{12}\selectfont\bfseries},
  subhead/.style={font=\fontsize{8.5}{10.2}\selectfont\bfseries,align=center},
  note/.style={font=\fontsize{7.6}{9.2}\selectfont,text=muted,align=center},
  visible/.style={circle,draw=ink,line width=.55pt,fill=visiblefill,minimum size=.40cm,inner sep=0pt},
  hidden/.style={circle,draw=edgeblue,line width=.7pt,fill=hiddenfill,minimum size=.45cm,inner sep=0pt},
  coupling/.style={draw=edgeblue,line width=.40pt,opacity=.43},
  spring/.style={draw=springgray,line width=.65pt,decorate,decoration={zigzag,segment length=2.0pt,amplitude=.65pt}},
  flow/.style={-{Stealth[length=2.0mm,width=1.35mm]},draw=ink,line width=.7pt},
  control/.style={-{Stealth[length=2.0mm,width=1.35mm]},draw=edgeblue,line width=.7pt},
  learnbox/.style={draw=rulegray,line width=.65pt,rounded corners=2pt,fill=visiblefill,minimum height=1.45cm,text width=2.2cm,align=center,inner sep=3pt}
}

% A 2 x 3 visible lattice; top-left coordinate and node-name prefix.
\newcommand{\visiblelattice}[3]{%
  \foreach \r in {0,1}{\foreach \c in {0,1,2}{
    \coordinate (#3p\r\c) at (#1+1.05*\c,#2-1.05*\r);
  }}
  \foreach \r in {0,1}{\foreach \c [evaluate=\c as \next using int(\c+1)] in {0,1}{
    \draw[spring] (#3p\r\c) -- (#3p\r\next);
  }}
  \foreach \c in {0,1,2}{\draw[spring] (#3p0\c) -- (#3p1\c);}
}
\newcommand{\visiblenodes}[1]{%
  \foreach \r in {0,1}{\foreach \c in {0,1,2}{\node[visible] at (#1p\r\c) {};}}
  \node[font=\fontsize{7.5}{8.5}\selectfont] at (#1p10) {$x_i$};
}

% Illustrative 8 x 8 thumbnail: blend a deterministic noise pattern and a
% hand-specified digit. Alpha is a drawing parameter, not simulated time/data.
\newcommand{\thumbnail}[3]{%
  \begin{scope}[shift={(#1,#2)},x=.22cm,y=.22cm]
    \foreach \value [count=\k from 0] in {
      0,0,0,0,0,0,0,0,
      0,0,85,100,100,65,0,0,
      0,75,35,0,15,100,45,0,
      0,0,0,0,30,95,0,0,
      0,0,0,70,100,55,0,0,
      0,0,0,0,15,85,70,0,
      0,70,30,0,10,90,65,0,
      0,0,70,100,95,65,0,0}{
      \pgfmathtruncatemacro{\row}{floor(\k/8)}
      \pgfmathtruncatemacro{\col}{mod(\k,8)}
      \pgfmathtruncatemacro{\noise}{mod(37*\row+23*\col+11*\row*\col+17,101)}
      \pgfmathtruncatemacro{\shade}{round((1-#3)*\noise+#3*\value)}
      \fill[black!\shade] (\col,-\row) rectangle ++(1,-1);
    }
    \draw[rulegray,line width=.5pt] (0,0) rectangle (8,-8);
  \end{scope}
}

\begin{document}
\mbox{}
\begin{tikzpicture}[remember picture,overlay,shift={(current page.south west)},x=1cm,y=1cm]
  \path[use as bounding box] (0,0) rectangle (19,13.2);

  % (a) Alternative hidden architectures. All displayed hidden units connect
  % to all displayed visible units; no hidden-hidden couplings are present.
  \node[paneltitle] at (.50,12.75) {(a) Visible and global hidden units};
  \node[subhead] at (2.30,12.18) {Ising: $K$ binary units};
  \node[subhead] at (6.60,12.18) {Potts: one $M$-state unit};
  \visiblelattice{1.25}{10.13}{i}
  \visiblelattice{5.55}{10.13}{p}
  \coordinate (ih1) at (1.55,11.33);
  \coordinate (ihK) at (3.05,11.33);
  \coordinate (potts) at (6.60,11.33);
  \foreach \r in {0,1}{\foreach \c in {0,1,2}{
    \draw[coupling] (ih1) -- (ip\r\c);
    \draw[coupling] (ihK) -- (ip\r\c);
    \draw[coupling] (potts) -- (pp\r\c);
  }}
  \visiblenodes{i}
  \visiblenodes{p}
  \node[hidden] at (ih1) {$h_1$};
  \node[hidden] at (ihK) {$h_K$};
  \node[text=muted] at (2.30,11.33) {$\cdots$};
  \node[note] at (2.30,11.85) {$h_j\in\{-1,+1\}$};
  \node[draw=edgeblue,fill=hiddenfill,rounded corners=3pt,line width=.7pt,minimum width=2.0cm,minimum height=.62cm] at (potts) {};
  \foreach \xx in {6.00,6.60,7.20}{\node[circle,draw=edgeblue,line width=.5pt,minimum size=.25cm,inner sep=0pt,fill=white] at (\xx,11.33) {};}
  \node[circle,fill=edgeblue,minimum size=.25cm,inner sep=0pt] at (6.60,11.33) {};
  \node[note] at (6.60,11.85) {one state active};
  \node[note,text width=8.0cm] at (4.45,8.53) {Continuous visible oscillators; nearest-neighbour springs};
  \draw[draw=edgeblue,line width=.8pt] (2.05,8.03) -- (2.50,8.03);
  \node[note,anchor=west] at (2.62,8.03) {Global hidden--visible couplings};
  \node[note,font=\fontsize{7.1}{8.5}\selectfont] at (4.45,7.74) {Optional per-pixel Ising units are omitted.};

  % (b) Training does require current parameters for hidden-weight gradients.
  % Aggregate signals live at their associated hidden node/state.
  \node[paneltitle] at (9.35,12.75) {(b) Local training at each noise time $t$};
  \node[learnbox] (clamp) at (10.70,10.85) {Clamp noisy\\data $\mathbf{x}\sim p_t$};
  \node[learnbox] (measure) at (13.80,10.85) {Measure\\$x_i$, $\psi_i=\partial_iF$\\hidden statistics};
  \node[learnbox] (update) at (16.90,10.85) {Evaluate local\\parameter\\updates};
  \draw[flow] (clamp.east) -- (measure.west);
  \draw[flow] (measure.east) -- (update.west);
  \node[note] at (16.90,12.05) {current couplings};
  \draw[control] (16.90,11.85) -- (update.north);
  \node[note,text width=6.6cm] at (13.00,9.55) {$\mathbf{x}$ stays clamped while hidden statistics\\and node-level aggregate signals are evaluated.};
  \node[draw=rulegray,fill=controlfill,rounded corners=2pt,minimum height=.66cm,text width=8.0cm,align=center] (schedule) at (13.95,8.55) {Store the learned parameter schedule $\theta(t)$};
  \draw[control] (update.south) -- (16.90,8.90);
  \node[note,text width=8.3cm] at (13.95,8.03) {Fast hidden equilibration; $\psi_i$ is the mean clamp force.};

  \draw[rulegray,line width=.65pt] (.5,7.55) -- (18.5,7.55);

  % (c) Open-loop schedule replay. There are no state-to-controller arrows.
  % Time labels use the implemented forward horizon t_max = 5.
  \node[paneltitle] at (.50,7.10) {(c) Generation by open-loop parameter replay};
  \node[fill=controlfill,draw=rulegray,line width=.6pt,rounded corners=3pt,minimum width=17.7cm,minimum height=.68cm] at (9.50,6.25) {Prelearned protocol $\theta(5-s)$: the same schedule for every generated sample};
  \foreach \xx in {3.30,9.50,15.70}{\draw[control] (\xx,5.90) -- (\xx,5.35);}
  \node[note] at (3.30,5.10) {$s=0,\quad t=5$};
  \node[note] at (9.50,5.10) {$s=2.5,\quad t=2.5$};
  \node[note] at (15.70,5.10) {$s=5,\quad t=0$};
  \thumbnail{2.42}{4.74}{0}
  \thumbnail{8.62}{4.74}{.60}
  \thumbnail{14.82}{4.74}{1}
  \draw[flow] (4.75,3.86) -- (8.05,3.86);
  \draw[flow] (10.95,3.86) -- (14.25,3.86);
  \node[note] at (6.40,4.12) {state evolves};
  \node[note] at (12.60,4.12) {state evolves};
  \node[subhead] at (3.30,2.66) {Noise initialization};
  \node[subhead] at (9.50,2.66) {Intermediate state};
  \node[subhead] at (15.70,2.66) {Generated sample};
  \node[font=\fontsize{10}{12}\selectfont] at (9.50,1.79) {$\mathrm{d}\mathbf{x}=\big[\mathbf{x}-2\nabla F_{\theta(5-s)}(\mathbf{x})\big]\,\mathrm{d}s+\sqrt{2T}\,\mathrm{d}\mathbf{B}_s$};
  \node[note] at (9.50,1.20) {Effective visible dynamics after averaging the fast hidden units};
  \node[note,text width=18cm,font=\fontsize{7.1}{8.5}\selectfont] at (9.50,.49) {Schematic network subsets and illustrative image sequence; no measured or generated samples are shown.};
\end{tikzpicture}
\end{document}
